% Declaración de uso de inteligencia artificial exigida por el instructivo
% N° 17/2026. Es el texto que estaba en la Sección 1.4, movido tal cual a las
% páginas preliminares: va inmediatamente antes de la introducción y conserva
% la numeración romana del frontmatter.
\clearpage
\chapter*{Declaración de uso de inteligencia artificial}
\addcontentsline{toc}{chapter}{Declaración de uso de inteligencia artificial}

En la elaboración de esta memoria se utilizaron herramientas de inteligencia
artificial generativa, concretamente los modelos Claude de Anthropic (familias
Opus y Sonnet 4.x), empleados principalmente a través de la aplicación de
escritorio en su modo Cowork. Estas herramientas se usaron a lo largo de todo
el trabajo: en la revisión de literatura y la comparación de
plataformas, en la discusión de decisiones de arquitectura, 
en la generación y refactorización de código, en la redacción y
edición de este documento,  y en tareas de apoyo
como el control de versiones y la organización del proyecto. Todo el contenido
generado o asistido fue revisado, evaluado y validado por el autor,
quien asume plena responsabilidad sobre el trabajo presentado. La
inteligencia artificial se empleó como herramienta de apoyo y no se reconoce
como coautoría.

El modo Cowork se utilizó como un acelerador del trabajo: permitió delegar la
realización del trabajo, como escribir código, redactar y editar, ejecutar y
verificar, conservando el rol del autor como arquitecto de las decisiones.
Trabajar de esta forma permitió mantener la constancia durante todo el semestre,
compatibilizando con las labores del autor como desarrollador de aplicaciones
móviles. La integración de la herramienta estuvo acompañado de la verificación:
la relectura crítica de cada resultado, la comparación de las afirmaciones y
planeamientos contra el código fuente, y el uso de una memoria persistente entre
sesiones para mantener el contexto y evitar contradicciones. El criterio
principal fue no dar por válido ningún resultado por el solo hecho de que
``funciona'' a primera vista.

Los prompts siguieron unos patrones recurrentes. Para una tarea de inicio a fin,
se partía entregando el contexto (usualmente un resumen de lo trabajado en
sesiones anteriores), seguido de lo que había que hacer y de los puntos que
valía la pena discutir, pidiendo siempre alternativas para elegir la
implementación más adecuada. Se cerraba con algún criterio de validación y con
notas que precisaban la ejecución. A mitad de una tarea el patrón cambiaba: se
revisaba lo ejecutado y se respondía con un punteo extenso de dudas, cambios
deseados y aspectos por considerar, y solo se confirmaba el avance al paso
siguiente cuando el resultado era satisfactorio. La mayoría de veces, se
recordaba al modelo lo que debería tener en memoria, para asegurar un correcto
alineamiento de la tarea a realizar. Los prompts breves se usaban para algunos
casos puntuales cuya discusión convenía expandir, y al final de una sesión
cargada se solicitaban los cierres habituales: formateo, guardado en memoria,
mensaje de \textit{commit} si se requería y un prompt de traspaso que orientara
a la sesión siguiente.

A modo ilustrativo, dos ejemplos representativos del tipo de instrucción
entregada:

\begin{itemize}
    \item Un prompt que funcionó bien: ``Retoma el informe: primero una foto del
    estado del repositorio con los commits nuevos. Luego valida afirmación por
    afirmación contra el código, marcando cada una como verdadera, falsa o
    imprecisa con su archivo y línea. Entrega un reporte ordenado por criticidad
    antes de editar y espera mi confirmación. Reglas: 80 caracteres, prosa en
    presente y opciones ante cada alternativa.'' Funcionó porque aportaba
    contexto, descomponía la tarea, exigía validación y opciones, y fijaba
    restricciones explícitas.
    \item Un prompt que funcionó mal: una instrucción amplia y sin contexto como
    ``mejora el capítulo de implementación'', que al no acotar el objetivo ni
    los criterios de aceptación producía cambios genéricos o de alcance excesivo
    que luego exigían una corrección considerable. La lección es que los prompts
    pobres en contexto y sin criterios claros rendían resultados de bajo valor.
\end{itemize}

Surgió un momento interesante en la etapa final del trabajo, al estudiar una
prueba de concepto de escape de contenedores para comprender el funcionamiento
de este tipo de \textit{exploit} y respaldar el análisis de contención del
entorno. Al
tratar un tema sensible, las primeras consultas sobre una prueba de concepto
realista del CVE~\cite{cve-runc-2025, poc-cve-2025-31133} en cuestión fueron
rechazadas por las políticas de uso de la herramienta. Tras solicitar al equipo
de Anthropic un acceso ampliado y explicar su propósito educativo dentro de este
trabajo de ciberseguridad, la consulta pudo retomarse. A partir de ahí, el
modelo resultó de gran ayuda para seguir el \textit{exploit} paso a paso,
consiguiendo
compararlo con otros CVE relacionados, distinguiendo, por ejemplo, los que se
ejecutan desde el interior del contenedor de los que afectan directamente al
\textit{host}.

Si el trabajo se reiniciara, se cambiarían algunas formas del uso de la
inteligencia artificial. Se establecería antes una rutina de revisión y
verificación sistemática, para no depender de que un resultado ``parezca''
correcto. Además, se sería más disciplinado con la memoria persistente y los
traspasos entre sesiones, que resultaron clave para mantener la coherencia a lo
largo de un trabajo asistido por IA.
